\begin{center}\textls[120]{b) \textsc{Harriot}s Lehrsatz. Zwei algebraische Theoreme.}\end{center}
\noindent\hspace*{\parindent}Der 1828 erschienene 3. Band von \textsc{Crelle}s Journal wird mit einem
kleinen Aufsatz von \textsc{Gauss} eröffnet: \textit{Beweis eines algebraischen Lehrsatzes} (Werke
III, S. 65), der einen einfachen Beweis für die \textsc{Descartes}sche Zeichenregel
(auch \textsc{Harriot}scher Lehrsatz genannt) gibt\footnote{Siehe \textsc{Descartes}, \textit{Geometrie}, deutsch von L. \textsc{Schlesinger}, 1923, S. 73. Der \textsc{Gauss}sche Beweis beruht auf der Abzählung der Änderungen, die die Anzahl der Zeichenwechsel in der Reihe der Koeffizienten eines Polynoms erfährt, wenn man dieses mit dem Faktor \(x-a\) für \(a>0\) multipliziert.} und diese Regel zugleich dahin
verschärft, dass die Anzahl der positiven Wurzeln einer Gleichung mit reellen
Koeffizienten eine ebenso grosse oder um eine gerade Zahl kleiner ist, als die Anzahl
der Zeichenwechsel in der Reihe der Koeffizienten. In einem Briefe an \textsc{Bessel}
vom 30. Mai 1828 (Werke X, 1, S. 248, 249) sagt \textsc{Gauss}, er habe »diesen
Beweis des \textsc{Harriot}schen Lehrsatzes \mbox{.\,.\,.} seitdem er in Göttingen sei, sehr
häufig mündlich vorgetragen«\footnote{Vergl. die Mitteilung von M. A. \textsc{Stern} über die Nachschrift eines um 1823 von \textsc{Gauss} gehaltenen Privatissimums, die diesen Beweis enthält, Göttinger Nachrichten 1869, S. 330, abgedruckt Werke X, 2, Abhandl. 1, S. 70.}, und \textsc{Pfaff} berichtet in einem Briefe vom
20. Oktober 1824 (Werke X, 1, S. 249, 250), er habe »vor einigen Jahren«
einen Beweis des gedachten Satzes als von \textsc{Gauss} herrührend in der Nachschrift
einer Vorlesung des Jenaer Professors v. \textsc{Münchow} gefunden; »ob dies
in einem Zusammenhang mit dem Erscheinen eines Beweises von derselben
Grundidee im letzten Hefte des 2. Bandes von \textsc{Crelle}s Zeitschrift \mbox{.\,.\,.} steht«,
fährt \textsc{Gauss} in jenem Briefe an \textsc{Bessel} fort\footnote{Gemeint ist der Aufsatz: J. A. \textsc{Grunert}, \textit{Beweis des Harriotschen Satzes}, \textsc{Crelle}s Journal 2, 1828, S. 335—344.}, »muss ich auf sich beruhen
lassen«. Dass die Beschäftigung mit diesem Satze für \textsc{Gauss} nicht nur
den Charakter einer Episode hatte, sondern dass ihn das Problem der Isolierung
\HistoricalFoot{}{24\textsuperscript{*}}
\par\clearpage\SourceLeaf{273}\ArithmeticHead{188}{L}{L. SCHLESINGER, ÜBER GAUSS' FUNKTIONENTHEORETISCHE ARBEITEN.}
\noindent der Wurzeln einer algebraischen Gleichung sehr interessiert hat, werden wir
weiter unten zu sehen Gelegenheit haben.

Im Herbst 1828 nahm \textsc{Gauss} an der damals in Berlin tagenden Naturforscherversammlung
teil; bei dieser Gelegenheit lernte er \textsc{Wilhelm Weber}
kennen, der dann Herbst 1831 als Nachfolger von \textsc{Johann Tobias Mayer} als
Professor der Physik nach Göttingen kam. Befördert durch das freundschaftliche
Zusammenleben mit \textsc{Weber}, das nur 1843—1849 eine Unterbrechung
erfuhr, tritt in den nächsten Jahren bei \textsc{Gauss} das Interesse für physikalische
Fragen in den Vordergrund, besonders wird der Erdmagnetismus ein Hauptgegenstand
seiner theoretischen Untersuchungen und seiner praktischen Betätigung\footnote{Über die Anfänge von \textsc{Gauss}' Beschäftigung mit dem Erdmagnetismus siehe E. \textsc{Schering}, \textit{C. F. Gauss und die Erforschung des Erdmagnetismus}, Abhandl. der k. G. d. Wissensch. zu Göttingen, 1887, Mathem. Klasse, S. 1—79, \textsc{Schering}s Werke II, S. 253 ff.; siehe auch Cl. \textsc{Schaefer}, Werke XI, 2, Abhandl. 2, S. 72 ff.}.

In die Zwischenzeit fallen noch zwei bedeutsame Publikationen; 1829 die
\textit{Theoria figurae fluidorum etc.} (Werke V, S. 29, \textit{Anzeige} S. 287) und 1831 die
zweite Abhandlung über die biquadratischen Reste (Werke II, S. 93, \textit{Anzeige}
S. 169). Aus der \textit{Theoria figurae fluidorum}\footnote{Vergl. hierfür \textsc{Bolza}, Abhandl. 5 dieses Bandes, S. 20—50.} interessiert uns hier besonders
eine Reihe von Sätzen, die für die Lehre vom Potential in der Ebene von
Wichtigkeit sind, so z. B. die Umformung eines Doppelintegrals in ein Linienintegral,
also der sog. \textsc{Green}sche Satz für die Ebene; aus der Abhandlung
über die biquadratischen Reste und der zugehörigen \textit{Anzeige} haben wir bereits
oben in den Abschnitten III d) S. 53 ff. und VIII c) S. 177, 178 das für unsern
Gegenstand vorwiegend in Betracht Kommende hervorgehoben. Das Jahr 1832
bringt dann als erste Frucht der Arbeit am Erdmagnetismus die \textit{Intensitas vis
magneticae terrestris etc.} (Werke V, S. 79, \textit{Anzeige} S. 293).

In den Göttingischen gelehrten Anzeigen vom 25. Februar 1833 erscheint
die Besprechung von \textsc{Fourier}s \textit{Analyse des équations déterminées} (1831) (Werke
III, S. 119—121), die uns in den Gedankenkreis des \textsc{Harriot}schen Lehrsatzes
zurückversetzt. \textsc{Gauss} charakterisiert den Zweck von \textsc{Fourier}s Werk geradezu
dahin, für eine Gleichung »die Anzahl der reellen und imaginären Wurzeln
\mbox{.\,.\,.} auszumitteln, jede\footnote{Nicht, wie es in beiden Drucken heisst, »jene«.} in feste Grenzen einzuschliessen und die stufenweise
\par\clearpage\SourceLeaf{274}\ArithmeticHead{189}{R}{IXb. ALGEBRAISCHE THEOREME.}
\noindent zu beliebiger Schärfe zu führende Annäherung zu ihren Werten methodisch
anzuordnen.« Als Grundlage hierfür erscheint der eine Generalisierung der
\textsc{Descartes}schen Zeichenregel gebende, sogenannte \textsc{Fourier}sche Lehrsatz\footnote{Theorem (A) S. 27 der Originalausgabe, S. 23 der \textsc{Loewy}schen Übersetzung, \textsc{Ostwald}s Klassiker Nr. 127. Über einige Ungenauigkeiten in der \textsc{Gauss}schen Anzeige vergl. die Bemerkungen II. von \textsc{Loewy}, Werke X, 1, S. 133, 134.}.
Jedoch »hat \textsc{Fourier} in das Wesen und in die Berechnung der imaginären
Wurzeln sich gar nicht eingelassen, und es bleibt daher noch ein weites Feld
zur Bearbeitung übrig« (Werke III, S. 121). Dass \textsc{Gauss} sich alsbald dieser
Bearbeitung zugewandt hat, geht aus einem Briefe an \textsc{Schumacher} vom
2. April 1833 hervor, wo er schreibt (Werke X, 1, S. 130): »Recht sehr danke
ich Ihnen \mbox{.\,.\,.} für die gütige Mitteilung des liniierten Papiers. Ich finde es
sehr brauchbar \mbox{.\,.\,.} auch zu Zeichnungen, die sich auf rein mathematische
Sachen beziehen, wie z. B. in Rücksicht auf die imaginären Wurzeln der
Gleichungen, und gerade mit Gegenständen der letzten Art habe ich mich in
der letzten Zeit viel beschäftigt. \mbox{.\,.\,.\,.} Wollte ich jene Untersuchungen, die,
wenn ich sie einmal entwickle, nur eine mässige Anzahl Bogen betragen
dürfen, mit der Breite, wie \textsc{Fourier}s Buch geschrieben ist, vortragen, so würde
ich vielleicht \mbox{.\,.\,.} mehrere grosse Quartbände gebrauchen.« Von diesen Untersuchungen
sind uns leider nur wenige Fragmente erhalten; es sind dies namentlich
zwei Theoreme, die \textsc{Gauss} in einem, sonst nur astronomischen Arbeiten
gewidmeten Handbuch (21, Bg) aufgezeichnet hat. Das erste (abgedruckt
Werke III, S. 112 und VIII, S. 32) besagt: Sind \(a_1,a_2,\ldots,a_n\) die Wurzeln
der Gleichung \(f(x)=0,\ a'_1,a'_2,\ldots,a'_{n-1}\) die der derivierten \(f'(x)=0\), und
denkt man sich in den \(a_1,\ldots,a_n\) gleiche anziehende oder abstossende Massen,
die im umgekehrten Verhältnis der Entfernung wirken, so findet in \(a'_1,\ldots,a'_{n-1}\)
Gleichgewicht statt. Bekanntlich kann man diesen Satz auch so aussprechen\footnote{Siehe z. B. \textsc{Osgood}, \textit{Lehrbuch der Funktionentheorie} I, 1907, S. 176.},
dass die Derivierte \(f'(x)\) ausserhalb des kleinsten konvexen Polygons oder des
kleinsten Intervalls, das alle Wurzeln der Gleichung \(f(x)=0\) in sich fasst,
nicht verschwinden kann, und in dieser Form erscheint er als Verallgemeinerung
des für reelle Wurzeln geltenden sogenannten \textsc{Rolle}schen Satzes. Das
zweite Theorem, das in dem Handbuch im unmittelbaren Anschluss an das
erste aufgezeichnet ist (abgedruckt Werke X, 1, S. 128), leitet \textsc{Gauss} mit den
\par\clearpage\SourceLeaf{275}\ArithmeticHead{190}{L}{L. SCHLESINGER, ÜBER GAUSS' FUNKTIONENTHEORETISCHE ARBEITEN.}
\noindent Worten ein: »Eine ganz neue Ansicht der Gleichungen beruhet auf Folgendem«.
Der Satz selbst besagt, dass die \(x\)-Ebene in \(n\) Gebiete zerfällt, von der Beschaffenheit,
dass 1) ein jedes eine Wurzel der Gleichung \(f(x)=0\) enthält,
2) jeder »Scheidungslinie« eine Wurzel der derivierten Gleichung \(f'(x)=0\)
angehört und 3), wenn \(f(x)=T+Ui\) gesetzt wird, längs einer jeden Scheidungslinie
\(U=\mathrm{const}\) ist. Die einleitenden Worte kehren fast unverändert
in einem Briefe an \textsc{Schumacher} vom 20. Juni 1836 wieder, wo es heisst
(Werke X, 1, S. 131): »Bei Gelegenheit der Vorlesung, die ich halte, bin ich
veranlasst in diesen Tagen auf die Theorie der Gleichungen zurückzukommen,
der ich jetzt einen ganz neuen Gesichtspunkt abgewonnen habe«. Man wird
daraus schliessen dürfen, dass die Aufzeichnung im Handbuch 21, Bg im Juni
1836 gemacht ist\footnote{\textsc{Fricke} gibt Werke VIII, S. 32 ohne weitere Begründung für den ersten Satz 1846 an.}. Und vier Tage darauf, am 24. Juni 1836 (Werke X, 1,
S. 132) schreibt \textsc{Gauss} an denselben Adressaten: »Ich habe in den letzten Tagen
meine Ideen über die Gleichungen weiter verfolgt. \mbox{.\,.\,.} Man würde um die
Verschiedenheit der Fälle wirklich anschaulich zu machen, eine grosse Menge
von Gleichungen in concreto durch Kurven versinnlichen müssen; jede Kurve
müsste durch Punkte gezeichnet werden und die Bestimmung eines einzigen
Punktes erfordert schon langwierige Rechnungen. Sie sehen es wohl der
Figur 4 bei meiner ersten Schrift von 1799\footnote{Werke III, auf der Tafel hinter S. 30; die Figur gibt den Verlauf der Kurven \(T=0,\ U=0\) für \(X=T+Ui=x^4-2x^2+3x+10\), siehe die Fussnote auf S. 24, ebenda.} nicht an, wie viel Arbeit die
richtige Zeichnung dieser Kurven erfordert hat, und doch ist dies vergleichsweise
nur ein sehr einfacher Fall gegen viele, die hier betrachtet werden
müssten«. Ähnlich heisst es viele Jahre später in einem Briefe an \textsc{Moebius}
vom 13. August 1849 (Werke X, 1, S. 109): »Ich möchte Ihrem Nachdenken
noch \mbox{.\,.\,.} empfehlen, \mbox{.\,.\,.} die verschiedene Gestaltung und Kreuzung der Linien
\(T=0\) und \(U=0\) in meiner Schrift von 1799 nach Massgabe der Ordnung
der betreffenden Funktion. Ich meine, die Enumeration aller verschiedenen
Fälle für die Konfiguration der unendlichen Äste. Anderes damit \textls[120]{verwandtes}
hat mich vielfach beschäftigt, und ich wollte erst zu meiner neulich in der
Sozietät gehaltenen Vorlesung\footnote{Gemeint ist die Jubiläumsschrift, siehe den folgenden Abschnitt; zur Theorie der Kurven \(T=0,\ U=0\) vergl. G. \textsc{Loria}, \textit{Spezielle algebraische und transzendente Kurven} I, 1910, S. 439—453.} die Darstellung der Hauptmomente dieser
Untersuchung als dritten Teil bestimmen.« Man kann hiernach, das was
\par\clearpage\SourceLeaf{276}\ArithmeticHead{191}{R}{IXb. ALGEBRAISCHE THEOREME.}
\noindent \textsc{Gauss} beabsichtigt hat, sich etwa in folgender Weise vorstellen. Die in dem
zweiten Theorem auftretenden Scheidungslinien sind die Kurven \(U=\mathrm{const}\),
die durch die Wurzelpunkte \(a'_1,\ldots,a'_{n-1}\) hindurchgehen, und man hat also
(vergl. die Bemerkung I, Werke X, 1, S. 133) für \(f(x)=X\), in der \(X\)-Ebene
durch die Punkte \(f(a'_k)\) zur reellen Achse parallele Halbstrahlen zu legen und
die entsprechenden Kurven in der \(x\)-Ebene zu konstruieren. Diese begrenzen
die im zweiten Theorem genannten \(n\) Gebiete, in deren jedem für jeden Wert
von \(X\) eine und nur eine Wurzel der Gleichung \(f(x)=X\) liegt. \textsc{Gauss} vollzieht
also durch diese Zerlegung der Ebene die Separation der \(n\) verschiedenen
Zweige der \(n\)-wertigen Funktion \(x\) von \(X\), ähnlich wie er für die
Modulfunktion in jedem der die rechte \(t\)-Halbebene bedeckenden Kreisbogenvierecke
einen eindeutig bestimmten Zweig der unendlich vielwertigen inversen
Funktion lokalisiert hat. Sein Bestreben war nun offenbar darauf gerichtet,
die »Konfiguration« jener \(U=\mathrm{const}\)-Kurven »nach Massgabe der Ordnung
der Funktion \(f(x)\)« zu ergründen, und dadurch Einsicht zu gewinnen in die
Art, wie die Lage der Wurzeln der Gleichung \(f(x)=X\) sich mit \(X\) ändert\footnote{In den beiden Beispielen Werke X, 1, S. 129 denkt er sich das absolute Glied \(g\) von \(f(x)\) variabel. Wir bemerken, dass in diesen Beispielen die Bezeichnungen \(\mathrm{Mod.}\,f(x)\) und \(\mathrm{Arg.}\,f(x)\) für absoluten Betrag und Argument benutzt werden. Später, in einer aus dem Jahre 1840 stammenden Aufzeichnung (Werke VIII, S. 335 ff.), nennt \textsc{Gauss} die konjugierte einer komplexen Grösse, die Adjunkte, eine Grösse vom absoluten Betrag Eins, eine zyklische Zahl, allgemein, wenn \(a=re^{\varphi i},\ r=\mathrm{Mod.}\,a\) ist, \(e^{\varphi i}\) den in \(a\) enthaltenen zyklischen Faktor und \(\varphi\) die Clise von \(a\).}.
Eine gewisse erste Orientierung hierfür liefert das erste Theorem, indem es
lehrt, dass für jeden Wert von \(X\), das die Wurzeln von \(f(x)=X\) umspannende
kleinste konvexe Polygon die \textls[120]{festen} Werte \(a'_1,\ldots,a'_{n-1}\) einschliessen
muss.

Bemerkenswert ist, wie die hier durchgeführte geometrisch-funktionentheoretische
Betrachtung gleichsam eine Brücke schlägt zwischen der der Theorie
des logarithmischen Potentials angehörenden Formulierung des ersten Theorems
und den ältesten einschlägigen Untersuchungen der Inauguraldissertation. Die
nächste, nunmehr zu besprechende Untersuchung wird uns zeigen, wie tief
\textsc{Gauss} in den Zusammenhang zwischen logarithmischem Potential und der
Lehre von der konformen Abbildung eingedrungen ist. —

